% Historical source (removed by author request): references/_OLD_Übungsaufgaben/Übungsblatt 3 - Vektorechnung/Blatt3.tex:358
\begin{ExerciseBlock}
\ExerciseHeading{14}{Geradengleichung -  Abstand Punkt-Gerade}{vector_algebra_point_line_distance}
Eine Gerade g verlaufe durch die 2 Punkte $P_1(2,1,3)$ und $P_2(-1,4,0)$.
 \begin{enumerate}[label=\alph*)]
	\item Stellen Sie die Geradengleichung der Geraden $g$ in der Parameterform auf.
	\item Berechnen Sie, ob der Punkt $Q(5,4,3)$ auf der Geraden liegt. Falls nein, berechnen Sie den Abstand des Punktes $Q$ zur Geraden $g$.
	\item Bestimmen Sie eine Gerade, die senkrecht auf $g$ steht und durch den Ursprung geht.
\end{enumerate}

\end{ExerciseBlock}
\UseExerciseSolution{vector_algebra_point_line_distance}
